% diversity_reduction_walkthrough.tex
%
% Self-contained walk-through of the diversity-reduction theory and the
% experiments that test it. Designed to be readable by someone with no prior
% exposure to the project. Includes all relevant math (with intuitive
% glosses; full proofs live in the dedicated theory sections) and a complete
% description of the experiments, their numerical results, and what they
% mean.
%
% This file is written as a single \section with subsections so it can be
% \input{} into a main paper as a stand-alone chapter, OR be promoted to a
% top-level chapter that supersedes the more compact diversity_reduction_*
% sections. After integration, the author can decide which proofs to migrate
% to an appendix.
%
% === INTEGRATION NOTES ===
%
% Cross-references to other paper sections are written defensively: each
% \ref{...} is paired with a parenthetical context phrase so the document
% remains readable even if the referenced section is moved or renumbered.
%
% Required packages: amsmath, amssymb, amsthm, booktabs, hyperref
% Required environments: theorem, proposition, corollary, lemma, definition,
%                        remark, proof (standard with amsthm)
%
% Citation keys assumed to exist in the bibliography:
%   turntrout2023steering, turntrout2023norms, brody2023expressivity,
%   sinii2025small, mardia2000directional
% (Same keys as diversity_reduction_theory.tex / _geometric.tex.)
%
% Code, configs, and per-run outputs for every empirical number reported here
% live in the project repository:
%   https://github.com/AMindToThink/steering_diversity
% under src/bounds/, scripts/bounds/, and outputs/bounds/.
%
% === END NOTES ===


\section{Activation Steering and Diversity: A Walk-Through}
\label{sec:walkthrough}

This section is a self-contained guide to the theory of steering-induced
diversity reduction and the experiments that test it. The reader is assumed
to know what a transformer is and what the residual stream is, but nothing
about steering vectors or the bounds in this paper. We develop the
mathematical picture from scratch, state and explain each claim, and then
walk through the experimental setup and results. Detailed proofs and a
small number of technical refinements live in the more compact
companion sections (Sections~\ref{sec:diversity-reduction},
\ref{sec:geometric-bound}, \ref{sec:conditions}, \ref{sec:synthesis}); we
refer to them inline.

\subsection{The Question}
\label{sec:walk-question}

\emph{Activation steering} is the practice of adding a fixed vector
$s \in \mathbb{R}^d$ to a transformer's residual stream at one or more
chosen layers, in order to bias the model's behaviour without retraining.
A steering vector trained to encode ``be more honest,'' ``be happier,''
or ``be more creative'' is added once per forward pass, then propagates
through the rest of the stack via the residual connections \citep{turntrout2023steering}.

Steering is cheap, training-free, and surprisingly effective at moving a
model's distribution of outputs in a desired direction. But practitioners
have repeatedly observed that as the steering strength is turned up, the
model's outputs become not just biased but \emph{narrow}: the same phrasings,
the same plot beats, the same fixed answer. Diversity collapses.

The question this paper answers is: \emph{why?} Specifically, we ask
whether the diversity collapse is a property of the architecture rather
than of any particular training procedure or steering vector. We show that
it is. Every transformer has normalization layers (LayerNorm or RMSNorm)
in every block, and these layers project the residual stream onto a sphere
of fixed radius. Adding a constant offset before this projection
\emph{must} concentrate the resulting representations toward a particular
``pole'' on that sphere, with the concentration tightening as the steering
magnitude grows. The mechanism is geometric, not behavioural.

The remainder of this section makes that intuition rigorous, derives
nine concrete bounds (some exact, some asymptotic), and tests all nine
on two open-weight models with three trained steering vectors and two
random-direction controls.

\subsection{Setup and Notation}
\label{sec:walk-setup}

Fix a transformer with hidden dimension $d$. Let $x \in \mathbb{R}^d$
denote the residual stream activation at a single token position of a
single prompt, drawn from some natural-language distribution $\mathcal{D}$.
We write
\[
  \mu \coloneqq \mathbb{E}_{x \sim \mathcal{D}}[x], \qquad
  \Sigma_x \coloneqq \mathbb{E}\!\left[(x - \mu)(x - \mu)^\top\right],
\]
for the unsteered mean and covariance. A steering intervention adds a
fixed vector $s \in \mathbb{R}^d$ to the residual stream, producing
$x' = x + s$, which is then passed through the next normalization step.

We focus on RMSNorm, which is what modern open-weight models (Llama,
Qwen, Mistral, etc.) use. RMSNorm is, up to a learned diagonal rescaling,
just \emph{projection onto a sphere}: it sends $v \mapsto v/\|v\|$
followed by multiplication by $\sqrt{d}$ and an elementwise gain vector
$\gamma$. Because the gain $\gamma$ is data-independent and the radius
$\sqrt{d}$ is a constant scale, the rank and angular structure of the
output are determined entirely by the bare projection
\begin{equation}
  g(v) \coloneqq \frac{v}{\|v\|}, \qquad g \colon \mathbb{R}^d \setminus \{0\} \to S^{d-1},
  \label{eq:walk-projection}
\end{equation}
which is what we analyze throughout. For LayerNorm, the only difference
is an additional projection that removes the all-ones component; this is
a second rank-reducing projection that does not change the qualitative
picture (see~\citealp{brody2023expressivity}).

It will be convenient to introduce one further piece of notation: the
\emph{shifted mean}
\begin{equation}
  m \coloneqq \mu + s, \qquad \hat{m} \coloneqq m / \|m\|,
\end{equation}
and its perpendicular projector
\begin{equation}
  P_{\perp m} \coloneqq I_d - \hat{m}\hat{m}^\top,
\end{equation}
which annihilates the radial direction $\hat{m}$ and leaves the
$(d-1)$-dimensional tangent space alone.

We also use the shorthand $\phi_s(x) \coloneqq \widehat{x + s} = (x+s)/\|x+s\|$
for the unit-normalized steered sample. The pushforward distribution
$\phi_s(\mathcal{D})$ lives on the unit sphere $S^{d-1}$. Two scalar
summary statistics will be central:
\begin{itemize}
  \item the \textbf{mean resultant vector} $\bar{R} \coloneqq \mathbb{E}[\phi_s(x)]$, of length $\|\bar{R}\| \in [0,1]$;
  \item the \textbf{spherical variance} $V \coloneqq 1 - \|\bar{R}\|$ \citep{mardia2000directional}.
\end{itemize}
$V = 0$ means the post-normalization distribution is a delta at one point on
the sphere; $V$ close to $1$ means the distribution is essentially uniform.

\subsection{The Geometric Picture}
\label{sec:walk-geometric}

The simplest and most robust statements come from a single elementary
lemma about how far apart the projections of two vectors can be.

\begin{lemma}[Chordal contraction]
\label{lem:walk-chordal}
For any nonzero $u, v \in \mathbb{R}^d$,
\begin{equation}
  \|\hat{u} - \hat{v}\|
  \;\le\;
  \frac{2\,\|u - v\|}{\max(\|u\|, \|v\|)}.
  \label{eq:walk-chordal}
\end{equation}
\end{lemma}

The proof is a one-line decomposition followed by the reverse triangle
inequality (full version in Section~\ref{sec:geometric-bound}, Lemma in
context). The takeaway: \emph{the chordal distance between two normalized
points is at most twice the Euclidean distance between the originals,
divided by their length}. Long vectors that differ by a fixed amount
project to nearby points on the sphere; short vectors that differ by the
same amount can project anywhere.

Apply this with $u = x + s$ and $v = s$ to compare a steered point to
the steering direction itself:

\begin{proposition}[Concentration around the pole; \textbf{Claim 3}]
\label{prop:walk-pole}
For any $x$ and any nonzero $s$,
\begin{equation}
  \|\phi_s(x) - \hat{s}\|
  \;\le\;
  \frac{2\,\|x\|}{\|s\|}.
  \label{eq:walk-pole}
\end{equation}
\end{proposition}

\begin{proof}
Setting $u = x+s$, $v = s$ in Lemma~\ref{lem:walk-chordal} gives
$\|\hat{u} - \hat{v}\| \le 2\|u-v\|/\max(\|u\|,\|v\|) = 2\|x\|/\max(\|x+s\|,\|s\|) \le 2\|x\|/\|s\|$.
\end{proof}

This is the geometric core of the whole story: \emph{every} post-steering
representation lies in a spherical cap around the pole $\hat{s}$, and the
cap shrinks linearly in $1/\|s\|$. The bound is exact, holds for every $x$
and every $\mathcal{D}$, and requires no distributional assumption.

The same lemma, applied with $u = x_1 + s$ and $v = x_2 + s$ (where
$s$ cancels in the numerator), gives a Lipschitz-style pairwise bound:

\begin{proposition}[Pairwise bound; \textbf{Claim 4}]
\label{prop:walk-pair}
If $\|s\| > R \coloneqq \max(\|x_1\|, \|x_2\|)$, then
\begin{equation}
  \|\phi_s(x_1) - \phi_s(x_2)\|
  \;\le\;
  \frac{2\,\|x_1 - x_2\|}{\|s\| - R}.
  \label{eq:walk-pair}
\end{equation}
\end{proposition}

The numerator $\|x_1 - x_2\|$ is the unsteered Euclidean distance: the
steering vector cancels. The denominator grows linearly in $\|s\|$, so
$\phi_s$ becomes a strict contraction on the unsteered support whenever
$\|s\|$ exceeds $R$ by enough. The catch is the precondition $\|s\| > R$:
if $\|s\|$ does not exceed the largest sample norm, the bound is vacuous
(the denominator is non-positive). We will see in
Section~\ref{sec:walk-results} that for realistic steering magnitudes at
the final-RMSNorm site this precondition is essentially never satisfied.

Three distribution-level corollaries follow immediately from
Proposition~\ref{prop:walk-pole}, by taking expectations and applying
the triangle inequality.

\begin{corollary}[Spherical variance bound; \textbf{Claim 2}]
\label{cor:walk-V}
\begin{equation}
  V(\phi_s(\mathcal{D}))
  \;\le\;
  \frac{2\,\mathbb{E}[\|x\|]}{\|s\|}.
  \label{eq:walk-V}
\end{equation}
\end{corollary}

\begin{corollary}[Expected pairwise distance; \textbf{Claim 6}]
\label{cor:walk-pair}
\begin{equation}
  \mathbb{E}\!\left[\|\phi_s(x_1) - \phi_s(x_2)\|\right]
  \;\le\;
  \frac{4\,\mathbb{E}[\|x\|]}{\|s\|}.
  \label{eq:walk-pair-exp}
\end{equation}
\end{corollary}

\begin{corollary}[Spherical diameter; \textbf{Claim 5}]
\label{cor:walk-diam}
For $\mathcal{D}$ supported on the ball of radius $R$,
\begin{equation}
  \operatorname{diam}\!\left(\phi_s(\mathcal{D})\right)
  \;\le\;
  \frac{4R}{\|s\|}.
  \label{eq:walk-diam}
\end{equation}
\end{corollary}

All four results (Claims 2, 3, 5, 6) come from the same one-line lemma.
They are exact upper bounds, distribution-free up to bounded support, and
all scale as $O(1/\|s\|)$.

\subsection{The Local Picture: Linearizing Around the Shifted Mean}
\label{sec:walk-local}

The geometric bounds are robust and exact, but they are scalar bounds:
they tell us how far things can be on the sphere, not \emph{which directions}
of variation are compressed. For that we need a local analysis, obtained
by Taylor-expanding the projection $g$ around $m = \mu + s$. Write
$x + s = m + \delta$ where $\delta = x - \mu$ has $\mathbb{E}[\delta] = 0$
and $\operatorname{Cov}(\delta) = \Sigma_x$.

A direct computation (Lemma~\ref{lem:jacobian} in
Section~\ref{sec:diversity-reduction}) gives the Jacobian of the
projection at $m$:
\begin{equation}
  J_g(m) = \frac{1}{\|m\|}\,P_{\perp m}.
  \label{eq:walk-jacobian}
\end{equation}
The Jacobian is precisely a rescaling by $1/\|m\|$ followed by projection
onto the tangent space at $\hat{m}$. Two effects pop out:
\begin{itemize}
  \item $P_{\perp m}$ \emph{annihilates the radial direction} $\hat{m}$.
  \item Everything that survives is \emph{rescaled by $1/\|m\|$}.
\end{itemize}
First-order in $\delta$, the Taylor expansion gives
\begin{equation}
  z = g(m + \delta) \approx \hat{m} + \frac{1}{\|m\|}\,P_{\perp m}\,\delta,
  \label{eq:walk-taylor1}
\end{equation}
so the steered output has mean approximately $\hat{m}$ and fluctuation
$P_{\perp m}\delta / \|m\|$. The output covariance is therefore
\begin{equation}
  \Sigma_z \approx \frac{1}{\|m\|^2}\,P_{\perp m}\,\Sigma_x\,P_{\perp m}.
  \label{eq:walk-output-cov}
\end{equation}

Two consequences fall out by direct algebra (Proposition~\ref{prop:diversity-reduction} in Section~\ref{sec:diversity-reduction}):

\begin{proposition}[Local covariance under steering; \textbf{Claim 1}]
\label{prop:walk-local}
Under the first-order approximation \eqref{eq:walk-taylor1},
\begin{enumerate}
  \item \textup{(Rank reduction)}
    $\operatorname{rank}(\Sigma_z) \le \min(\operatorname{rank}(\Sigma_x), d-1)$.
    If $\hat{m}^\top \Sigma_x \hat{m} > 0$, then $\operatorname{rank}(\Sigma_z) \le \operatorname{rank}(\Sigma_x) - 1$.
  \item \textup{(Total variance shrinks as $1/\|m\|^2$)}
    $\operatorname{tr}(\Sigma_z) = \dfrac{\operatorname{tr}(\Sigma_x) - \hat{m}^\top \Sigma_x \hat{m}}{\|m\|^2}$.
\end{enumerate}
\end{proposition}

The numerator is the unsteered total variance \emph{minus the variance
along the steered direction}; this is non-negative because $\hat{m}^\top \Sigma_x \hat{m}$
is bounded above by the largest eigenvalue of $\Sigma_x$. The denominator
grows as $\|m\|^2$, which in the strong-steering regime ($\|s\| \gg \|\mu\|$)
is approximately $\|s\|^2$. So the local bound predicts
$\operatorname{tr}(\Sigma_z) = O(1/\|s\|^2)$.

This is consistent with the global bound's $O(1/\|s\|)$ on chordal distance
because, for distributions concentrated near a pole, chordal distance
$\theta$ scales linearly with angular displacement and variance scales as
$\theta^2$. Squaring $O(1/\|s\|)$ gives $O(1/\|s\|^2)$. The two analyses
are not in tension; they are dual lenses on the same phenomenon.

\subsection{Second-Order Refinement: Where the Parabola Comes From}
\label{sec:walk-second-order}

The first-order expansion treats the projection $g$ as if it were exactly
linear. It is not: the sphere is curved, and the correction matters when
the input variation $\delta$ is not negligibly small relative to $\|m\|$.
A second-order Taylor expansion (full derivation in
Section~\ref{sec:second-order} of the appendix) gives, after the
$\alpha^2$ terms cancel,
\begin{equation}
  z \;\approx\; \left(1 - \frac{\|\delta_\perp\|^2}{2\|m\|^2}\right)\hat{m}
  \;+\; \left(\frac{1}{\|m\|} - \frac{\alpha}{\|m\|^2}\right)\delta_\perp,
  \label{eq:walk-taylor2}
\end{equation}
where $\alpha = \hat{m}^\top \delta$ is the radial component of $\delta$
and $\delta_\perp = P_{\perp m}\,\delta$ is the perpendicular component.
Two effects appear:
\begin{enumerate}
  \item \emph{Sphere curvature} pulls the radial coefficient below $1$
    by $\|\delta_\perp\|^2 / (2\|m\|^2)$ --- the perpendicular spread
    causes the mean resultant vector to retreat from the surface of the
    sphere. This is the \emph{quadratic} term in the variance.
  \item \emph{Radial-perpendicular coupling}: if a sample is further from
    the origin ($\alpha > 0$), its angular displacement is suppressed;
    if it is closer ($\alpha < 0$), it is amplified. This is the
    perturbative manifestation of the magnifying-glass effect we will
    discuss in Section~\ref{sec:walk-conditions}.
\end{enumerate}
Taking expectations yields
\begin{equation}
  V \;\approx\; \frac{\operatorname{tr}(\Sigma_x) - \hat{m}^\top \Sigma_x \hat{m}}{2\|m\|^2}.
  \label{eq:walk-V-second-order}
\end{equation}
This refines Corollary~\ref{cor:walk-V}: the spherical variance is not
just $\le 2\mathbb{E}\|x\|/\|s\|$, it is approximately equal to a specific
quadratic form divided by $2\|m\|^2$. Substituting $m = \mu + s$ gives
\begin{equation}
  \|m\|^2 \;=\; s^2 \|v_{\mathrm{eff}}\|^2 \;+\; 2 s\, (v_{\mathrm{eff}} \cdot \mu) \;+\; \|\mu\|^2,
  \label{eq:walk-m-quadratic}
\end{equation}
where $v_{\mathrm{eff}}$ is the per-unit-scale steering direction. So
$\|m(s)\|^2$ is a \emph{quadratic} in $s$, with floor $\|\mu\|^2$ at
$s = 0$. On a log--log plot of $V$ vs $s$, this denominator structure
produces a smooth crossover: a flat shelf at $s \ll \|\mu\|/\|v_{\mathrm{eff}}\|$,
a slope $-2$ tail at $s \gg \|\mu\|/\|v_{\mathrm{eff}}\|$, and a
``parabolic-looking'' bend in between. We will return to this in
Section~\ref{sec:walk-second-order-check}, where it lets us \emph{recover}
the unsteered mean norm $\|\mu\|$ from the curvature of $V(s)$.

\subsection{When Does Steering Actually Reduce Diversity?}
\label{sec:walk-conditions}

The bounds above tell us how concentrated the post-steering distribution
is on the sphere. They do not, by themselves, say whether this is
\emph{less} concentrated than the unsteered distribution. For that we
need to compare $\|m\|$ to $\|\mu\|$. The first-order trace formula is
$\operatorname{tr}(\Sigma_z) \propto 1/\|m\|^2$, so steering reduces
diversity (in this sense) when
\begin{equation}
  \|\mu + s\| > \|\mu\|.
  \label{eq:walk-condition}
\end{equation}
Squaring and rearranging:
\begin{equation}
  2\,\mu^\top s + \|s\|^2 > 0,
  \qquad\text{equivalently}\qquad
  \cos\angle(\mu, s) > -\frac{\|s\|}{2\|\mu\|}.
\end{equation}
This is a very weak condition. The right-hand side is always negative,
so the condition is automatic whenever $s$ is in the same hemisphere as
$\mu$, and even tolerates obtuse angles up to
$\arccos(-\|s\|/2\|\mu\|)$. Only an $s$ that is both nearly anti-parallel
to $\mu$ \emph{and} small relative to $\|\mu\|$ can violate it.

When the condition is violated, $\|m\|$ is small and the projection
becomes a magnifying glass: small Euclidean differences become large
angular differences, and the diversity \emph{increases} rather than
decreases. This anti-pole regime is analyzed in
Section~\ref{sec:conditions} (companion to this walkthrough).

In high dimensions, $\cos\angle(\mu, s)$ tends to be small in absolute
value for any pair of structured vectors that are not deliberately aligned;
combined with the typical residual-stream norm growth across layers
\citep{turntrout2023norms}, the condition is generically satisfied. But
``generically'' is not ``always'' --- one of our experiments
(Section~\ref{sec:walk-claim7}) lands on the wrong side of the condition
in a way that has practical consequences for how steering vectors are
deployed. We highlight it as the most important
``check before applying'' caveat in the entire bounds story.

\subsection{Asymptotic Scaling Law and Eigenstructure}
\label{sec:walk-scaling}

Two final predictions are corollaries of the local analysis combined with
the relationship between chordal and angular spread on the sphere
(Section~\ref{sec:synthesis}):

\begin{enumerate}
  \item \textbf{Scaling law (Claim 8).} As $\|s\| \to \infty$,
    $V \sim 1/\|s\|$ and $\operatorname{tr}(\Sigma_z) \sim 1/\|s\|^2$.
    On a log--log plot, slopes of $-1$ and $-2$ respectively.
  \item \textbf{Alignment (Claim 9).} The smallest-eigenvalue eigenvector
    of $\Sigma_z$ aligns with $\hat{m}$, since $\Sigma_z$'s null direction
    \emph{is} $\hat{m}$ in the first-order picture.
\end{enumerate}

These are the most direct empirical predictions: a single regression on
the log--log plot tests Claim 8, and a single eigendecomposition tests
Claim 9. We test both in Section~\ref{sec:walk-results}.

\subsection{Summary of the Nine Claims}

For convenience, Table~\ref{tab:walk-claims} restates the nine bounds in
one place. Numbers refer to the order in which they were originally derived
in the project; we also indicate whether the claim is exact (a strict
inequality from Lemma~\ref{lem:walk-chordal}) or asymptotic (relies on
truncating a Taylor expansion).

\begin{table}[htbp]
  \centering
  \small
  \caption{The nine claims tested empirically.}
  \label{tab:walk-claims}
  \begin{tabular}{@{}clp{0.55\linewidth}l@{}}
    \toprule
    \# & Quantity bounded & Bound & Type \\
    \midrule
    1 & $\operatorname{tr}(\Sigma_z)$ & $\approx (\operatorname{tr}\Sigma_x - \hat{m}^\top \Sigma_x \hat{m}) / \|m\|^2$ & Taylor 1st order \\
    2 & $V$ & $\le 2\,\mathbb{E}\|x\| / \|s\|$ & exact \\
    3 & $\max_x \|\phi_s(x) - \hat{s}\|$ & $\le 2\|x\| / \|s\|$ & exact \\
    4 & $\|\phi_s(x_1) - \phi_s(x_2)\|$ & $\le 2\|x_1-x_2\| / (\|s\|-R)$, $\|s\|>R$ & exact \\
    5 & $\operatorname{diam} \phi_s(\mathcal{D})$ & $\le 4R / \|s\|$ & exact \\
    6 & $\mathbb{E}\|\phi_s(x_1) - \phi_s(x_2)\|$ & $\le 4\,\mathbb{E}\|x\| / \|s\|$ & exact \\
    7 & ``Reduction gate'' & $\|m+s\| > \|m\|$ ($\Leftrightarrow \cos > -\|s\|/2\|\mu\|$) & algebraic \\
    8 & Scaling law & $V \sim 1/\|s\|$, $\operatorname{tr}\Sigma_z \sim 1/\|s\|^2$ as $\|s\| \to \infty$ & asymptotic \\
    9 & Eigenalignment & smallest eigvec of $\Sigma_z$ $\parallel$ $\hat{m}$ & asymptotic \\
    \bottomrule
  \end{tabular}
\end{table}

\subsection{Experimental Setup}
\label{sec:walk-setup-empirical}

\paragraph{Models.} We use two open-weight checkpoints that span an order
of magnitude in size:
\begin{itemize}
  \item Qwen2.5-1.5B-Instruct: $d = 1536$, $L = 28$ blocks, RMSNorm + RoPE.
  \item Meta-Llama-3-8B-Instruct: $d = 4096$, $L = 32$ blocks, RMSNorm + RoPE.
\end{itemize}
RoPE applies positional information \emph{inside attention} (rotating Q
and K), not additively to the residual stream, so the residual stream
itself never carries a positional embedding. None of our quantities are
contaminated by a positional bias.

\paragraph{Data.} Each run streams $1000$ documents from
\texttt{HuggingFaceFW/fineweb-edu} (a quality-filtered web-text corpus),
truncates each to $256$ tokens, and runs one forward pass per document.
Token-level activations (with padding masked) are the unit of statistics.
We deliberately use a broad input distribution rather than a domain-matched
one, to stress rather than cherry-pick the bounds.

\paragraph{Steering vectors.} We test three publicly released vectors
from the EasySteer
library\footnote{\url{https://github.com/ZJU-REAL/EasySteer}}, plus two
random controls:
\begin{itemize}
  \item \texttt{happy\_diffmean.gguf} (Qwen): happy-vs-sad diff-mean from role-play prompts.
  \item \texttt{style-probe.gguf} (Qwen): a stylistic-contrast direction.
  \item \texttt{create.gguf} (Llama): a creativity direction trained on Llama-3-8B.
  \item \textbf{Per-layer-matched random}: at each target layer $\ell$, draw
    $r_\ell \sim \mathcal{N}(0, I_d)$ and rescale so $\|r_\ell\| = \|s_\ell\|$.
  \item \textbf{Aggregate-matched random}: same draws, but rescale all
    per-layer vectors by a single common factor so that
    $\|\sum_\ell r_\ell\| = \|\sum_\ell s_\ell\|$ over the target layers.
\end{itemize}
The two random controls differ in a way that turns out to be central
(Section~\ref{sec:walk-agg-matched}). Random unit vectors in high
dimensions are near-orthogonal, so $\|\sum r_\ell\|$ scales as $\sqrt{n}$;
trained vectors stack coherently so $\|\sum s_\ell\|$ scales as $n$.
Per-layer matching at the same nominal scale therefore delivers
\emph{less} aggregate steering for random than for trained.

\paragraph{Pipeline.} For each (model, vector, target layers, random-mode)
configuration and each nominal scale $c \in \{0, 0.5, 1, 2, 4, 8\}$, the
pipeline:
\begin{enumerate}
  \item runs $1000$ documents through the model in bfloat16,
    intervening at the chosen target layers using \texttt{nnsight};
  \item captures the pre- and post-RMSNorm residual at the final
    RMSNorm site, masking padding tokens;
  \item streams a full $d \times d$ covariance $\Sigma_x$ via a batched
    Chan--Golub--LeVeque Welford update (in float32, validated against
    \texttt{numpy.cov} as ground truth);
  \item maintains a reservoir of $K = 1024$ representative activation
    vectors (Vitter Algorithm R) for pairwise statistics;
  \item produces \texttt{stats.pt} and \texttt{stats\_meta.json}, which
    feed the per-claim verification step.
\end{enumerate}
A hard ``trust no silent steering'' verification step runs before any
recording: the verifier samples from the model with and without steering
and refuses to record statistics unless the steered samples differ from
the unsteered ones on both KL-divergence and residual-delta-magnitude
criteria. This gate caught the batch-indexing bug discussed in
Section~\ref{sec:walk-threats}.

\paragraph{Two experiments.} The full protocol is $19$ recording runs:
\begin{itemize}
  \item \textbf{Experiment 1} ($5$ runs): Qwen + \{happy, style, per-layer
    random matched to happy\}; Llama + \{creativity, per-layer random
    matched to creativity\}. Target layers are the full trained range
    for each real vector.
  \item \textbf{Experiment 2} ($14$ runs): two aggregate-matched random
    controls (one per model), plus a single-layer sweep at edge/middle/edge
    of each trained range (Qwen happy at layers $\{10, 17, 25\}$; Llama
    creativity at layers $\{16, 22, 29\}$) for both real and per-layer
    random vectors. This experiment is what teases apart aggregate
    magnitude from semantic direction.
\end{itemize}

\paragraph{Effective steering.} The theory above treats $s$ as if it were
added directly at the measurement layer. In practice $s$ is added at
target layers (typically middle layers), and the residual then propagates
through additional blocks before reaching the final RMSNorm. Writing the
theoretical RHS in terms of the residual at the measurement layer
requires an ``effective steering'' quantity. We use
\[
  s_{\mathrm{eff}} \coloneqq \mu_c - \mu_0,
\]
the empirical mean shift of the pre-RMSNorm residual between the steered
and unsteered runs at the same scale. By construction $\|s_{\mathrm{eff}}\|$
is monotone in $c$. All theoretical RHS quantities are reported using
$s_{\mathrm{eff}}$ in place of $s$.

\subsection{Results, Claim by Claim}
\label{sec:walk-results}

Table~\ref{tab:walk-summary} collects the headline result for every claim.
The remainder of this section discusses the noteworthy ones in turn. Per-run
plots of every LHS-vs-RHS curve are in
\texttt{outputs/bounds/<run\_name>/plots/} in the repository.

\begin{table}[htbp]
  \centering
  \small
  \caption{Summary of empirical findings, aggregated across all $19$ runs.}
  \label{tab:walk-summary}
  \begin{tabular}{@{}llp{0.55\linewidth}@{}}
    \toprule
    \# & Status & Detail \\
    \midrule
    1 & qualitative & Empirical ratio LHS/RHS $\approx 0.1$--$0.5$; right direction, $2$--$10\times$ loose. \\
    2 & holds & Passes at every scale on every run. Tightness varies. \\
    3 & holds & Passes everywhere; strict geometric identity. \\
    4 & inapplicable & Precondition $\|s_{\mathrm{eff}}\| > R$ never met at the final-RMSNorm site. \\
    5 & holds & Passes everywhere. \\
    6 & holds & Passes everywhere; typically loose. \\
    7 & \textbf{fails on Qwen + happy} & At every scale; also fails for Qwen single-layer at $L17, L25$ (direction-independent). \\
    8 & \textbf{magnitude not direction} & Manifests on Llama; fails to manifest on Qwen at any tested scale. \emph{Aggregate-matched random matches trained}. \\
    9 & not observed & $|\cos|$ between smallest eigenvector of $\Sigma_z$ and $\hat{m}$ stays $\approx 0$ across all runs. \\
    \bottomrule
  \end{tabular}
\end{table}

\subsubsection{The exact bounds (Claims 3, 5)}

Proposition~\ref{prop:walk-pole} (concentration around the pole) and
Corollary~\ref{cor:walk-diam} (spherical diameter) are strict inequalities
derived from Lemma~\ref{lem:walk-chordal}. They hold for every $x$, every
$s$, and every distribution. As expected, they pass at every scale in
every run and would only fail in the presence of a pipeline bug. We use
them as a smoke test before reporting any other claim.

\subsubsection{The Lipschitz bound (Claim 4) is never applicable}

The pairwise bound of Proposition~\ref{prop:walk-pair} requires
$\|s\| > R$. At the final-RMSNorm site, $R$ (the largest sample norm in
the support) is on the order of hundreds for both models, while
$\|s_{\mathrm{eff}}\|$ at the maximum tested scale is around $120$. The
precondition is never satisfied, so we report Claim 4 as
\emph{always inapplicable} rather than reporting an LHS/RHS ratio with a
negative denominator.

This is not an indictment of the proposition --- it remains an exact
inequality whenever the precondition holds. It does suggest that the
right place to evaluate the pairwise bound is the steering-injection
layer, not the final RMSNorm.

\subsubsection{The first-order trace bound (Claim 1) is loose but right-direction}

Proposition~\ref{prop:walk-local} predicts
$\operatorname{tr}(\Sigma_z) \approx (\operatorname{tr}\Sigma_x - \hat{m}^\top \Sigma_x \hat{m}) / \|m\|^2$.
Empirically, the ratio LHS/RHS is in $[0.1, 0.5]$ across runs. The
prediction is loose because we are at a measurement site many blocks
downstream of the steering site; intermediate attention and MLP blocks
re-shape the distribution. The direction of the prediction is correct,
which is what we should expect from a first-order Taylor truncation; we
get a tighter, asymptotically-correct prediction from the second-order
analysis, which we test directly in Section~\ref{sec:walk-second-order-check}.

\subsubsection{The reduction gate (Claim 7) fails for Qwen + happy and at single layers}
\label{sec:walk-claim7}

The reduction gate of \eqref{eq:walk-condition} states
$\|\mu + s\| > \|\mu\|$. It is an algebraic identity and so cannot fail
``in theory''; it can fail \emph{in practice} when the empirical
$\mu \cdot s$ is just-negative-enough to dominate $\|s\|^2/2$ at small
$\|s\|$.

In Experiment 1, four of the five configurations pass at every scale.
One configuration --- Qwen + \texttt{happy} --- \textbf{fails at every
scale}. A diagnostic
(\texttt{scripts/bounds/investigations/03\_investigate\_happy\_data\_alignment.py})
revealed why: on FineWeb-Edu,
$\cos(\mu_{\mathrm{unsteered}}, \sum_\ell s^{\mathrm{happy}}_\ell) = -0.108$,
about $4\sigma$ from the null distribution of random-direction cosines in
$d = 1536$. Twelve of $16$ per-layer cosines are negative. The
\texttt{happy\_diffmean} vector was trained on a happy-vs-sad role-play
contrast, and educational prose at Qwen's representation lives slightly
on the sad side of that axis. The result is that $2\mu^\top s$ remains
just-large-enough negative to flip the sign of $\|m\|^2 - \|\mu\|^2$ at
all tested scales.

Experiment 2 produced a second, more surprising failure pattern. At Qwen
single-layer steering at \emph{layers $L17$ and $L25$}, Claim 7 fails for
both real and random directions. At single-layer $L10$, it passes for
both. This is direction-independent and position-dependent. The natural
interpretation is that early-layer steering has many downstream blocks to
amplify $\|s_{\mathrm{eff}}\|$, late-layer steering does not, and at small
effective magnitude the sign of $\mu \cdot s_{\mathrm{eff}}$ decides.

The practical lesson is sharp: the reduction gate is a \emph{precondition}
that should be checked empirically before applying any of the
diversity-reduction bounds to a new $(\text{model}, \text{vector}, \text{dataset})$
triple.

\subsubsection{The scaling law (Claim 8) and the magnitude-vs-direction story}
\label{sec:walk-agg-matched}

Claim 8 predicts log--log slope $-1$ for $V(\|s\|)$ and slope $-2$ for
$\operatorname{tr}(\Sigma_z)(\|s\|)$ in the asymptotic regime. We
estimate slopes by log--log linear regression across $c \in \{0.5, 1, 2, 4, 8\}$
for every run.

\begin{table}[htbp]
  \centering
  \small
  \caption{Empirical slopes from Claim 8 regressions. Predictions: slope $V = -1$, slope $\operatorname{tr}(\Sigma_z) = -2$.}
  \label{tab:walk-scaling}
  \begin{tabular}{@{}lrr@{}}
    \toprule
    Run & slope $V$ & slope $\operatorname{tr}(\Sigma_z)$ \\
    \midrule
    \multicolumn{3}{l}{\emph{Llama-3-8B-Instruct:}} \\
    \texttt{llama\_creativity} (trained, full target range) & $\mathbf{-1.92}$ & $-1.77$ \\
    \texttt{llama\_random\_agg\_matched} (aggregate-matched random) & $\mathbf{-1.85}$ & $-1.71$ \\
    \texttt{llama\_random} (per-layer-matched random) & $-0.51$ & $-0.41$ \\
    \texttt{llama\_creativity\_single\_L\{16,22,29\}} & $\approx -0.05$ & $\approx -0.03$ \\
    \texttt{llama\_random\_single\_L\{16,22,29\}} & $\approx -0.04$ & $\approx -0.02$ \\
    \midrule
    \multicolumn{3}{l}{\emph{Qwen2.5-1.5B-Instruct:}} \\
    \texttt{qwen\_style} & $-0.29$ & $-0.25$ \\
    \texttt{qwen\_random\_agg\_matched} (aggregate-matched random) & $-0.20$ & $-0.17$ \\
    \texttt{qwen\_random} (per-layer-matched random) & $-0.04$ & $-0.03$ \\
    \texttt{qwen\_happy} & $-0.02$ & $-0.02$ \\
    \texttt{qwen\_*\_single\_L\{10,17,25\}} & $\approx 0$ & $\approx 0$ \\
    \bottomrule
  \end{tabular}
\end{table}

Two findings dominate Table~\ref{tab:walk-scaling}.

\paragraph{Llama hits the scaling regime; Qwen does not.}
On Llama, both the trained creativity vector and the aggregate-matched
random control hit slope-$V$ values steeper than the predicted $-1$
($-1.92$ and $-1.85$ respectively). The overshoot reflects the fact that
at $c = 4, 8$ Llama is already well into the regime where its outputs
collapse into gibberish, pushing the spherical distribution past where
the first-order asymptotic captures the behaviour.

On Qwen no run reaches the asymptotic regime. The reason is a magnitude
mismatch: Qwen's $\|\mu\|$ at the final RMSNorm site is approximately
$218$ and $\mathbb{E}\|x\| \approx 300$, whereas $\|s_{\mathrm{eff}}\|$ at
$c = 8$ is only about $125$. The bounds become tight when
$\|s_{\mathrm{eff}}\|$ overtakes the natural scale of the residual stream;
Qwen is an architecture for which this never happens at scales we can
push the model through without destroying its outputs.

\paragraph{Aggregate magnitude beats semantic direction.}
The cleanest finding of the paper is the side-by-side comparison on Llama:
trained creativity has slope $V = -1.92$; aggregate-matched random has
slope $V = -1.85$. The difference of $0.07$ is at the level of regression
noise. \textbf{Random directions, when given matched aggregate magnitude,
drive the scaling law just as hard as trained directions.}

This contradicts the obvious reading of Experiment 1 alone, where the
per-layer-matched random had slope $-0.51$ and the trained vector had
slope $-1.92$. That apparent ``trained vectors are special'' gap was an
artifact of how the per-layer-matched random under-delivers aggregate
steering: with $14$ near-orthogonal random unit vectors, the aggregate
norm $\|\sum_\ell r_\ell\|$ grows as $\sqrt{14} \approx 3.7$, whereas the
trained vectors stack coherently so $\|\sum_\ell s_\ell\|$ grows as $14$.
At the same nominal scale $c$, per-layer random delivers $\sim 4\times$
less aggregate steering.

The implication for the bounds: \emph{Claim 8 is a statement about
aggregate steering magnitude, not about semantic direction}. Any
direction with sufficient aggregate magnitude reproduces the scaling.
And the implication for empirical practice: random-direction baselines
should be \emph{aggregate-matched}, not per-layer-matched, on pain of
underestimating how hard the bound is being driven.

The single-layer sweeps tell a complementary story: with only one
intervention layer, the aggregate magnitude is much smaller, so slopes
collapse toward zero (single-layer Llama runs are all in
$[-0.06, -0.03]$, single-layer Qwen runs are indistinguishable from
zero). Real and random are indistinguishable at matched single-layer
magnitude.

\subsubsection{The eigenalignment claim (Claim 9) is not observed}

Across all $19$ runs, the absolute cosine between the smallest-eigenvalue
eigenvector of $\Sigma_z$ and $\hat{m}$ stays close to zero with no
monotone trend in $c$. Three explanations are plausible: (i) the
reservoir-based estimate of $\Sigma_z$ ($K = 1024$ samples in
$d = 1536$ or $4096$) is too noisy to pin down the smallest eigenvector;
(ii) the theoretical $\hat{m}$ is the steering direction at the
\emph{injection layer}, but at the final-RMSNorm measurement site the
relevant axis is a propagated and rotated version of this; (iii) the
prediction may simply not survive realistic distributional conditions.
We do not currently distinguish between these and recommend that future
work test the alignment claim at the injection layer rather than the
measurement layer.

\subsection{Second-Order Consistency Check: The Parabola Has a Meaning}
\label{sec:walk-second-order-check}

The log--log plots of $V(s)$ in our scale sweeps are not straight lines.
They are smooth curves --- ``parabola-like'' on log--log axes. The
second-order theory of Section~\ref{sec:walk-second-order} predicts
exactly this shape. From \eqref{eq:walk-V-second-order} and
\eqref{eq:walk-m-quadratic}:
\begin{equation}
  V(s) \;\approx\; \frac{\operatorname{tr}\Sigma_x - \hat{m}^\top \Sigma_x \hat{m}}{2\|m(s)\|^2},
  \qquad
  \|m(s)\|^2 = a s^2 + b s + c,
  \quad c = \|\mu\|^2.
  \label{eq:walk-V-parabola}
\end{equation}
This makes two independent, falsifiable predictions:
\begin{enumerate}
  \item $\|m(s)\|^2$ is a perfect quadratic in $s$ when expanded around
    the measurement-site residual. Fitting $a s^2 + b s + c$ to the
    measured $\|m(s)\|^2$ at $s > 0$ and \emph{extrapolating to $s = 0$}
    must recover $\sqrt{c} = \|\mu\|$ measured directly at $s = 0$.
  \item In the deep-asymptotic regime where the second-order Taylor
    expansion has converged, the constant
    $V \cdot \|m\|^2$ should equal $(\operatorname{tr}\Sigma_x - \hat{m}^\top \Sigma_x \hat{m})/2$.
\end{enumerate}

We test both directly. The script
\texttt{scripts/bounds/second\_order\_consistency\_check.py} extracts
$(s, V, \|m\|, \operatorname{tr}\Sigma_x, \hat{m}^\top \Sigma_x \hat{m})$
from each \texttt{stats.pt} file, fits the quadratic, and writes the table
below. The independently-measured $\|\mu\|$ is read off the raw
\texttt{mean} vector at $s = 0$ in the same \texttt{stats.pt}, with no
fitting involved.

\begin{table}[htbp]
  \centering
  \small
  \caption{Second-order consistency check: the quadratic intercept $\sqrt{c}$ recovers $\|\mu\|$ to within ~1--9\% across all runs where the model is valid in the small-$s$ regime.}
  \label{tab:walk-sqrtc}
  \begin{tabular}{@{}lrrrr@{}}
    \toprule
    Run & R$^2$ of $\|m\|^2$ fit & $\sqrt{c}$ & $\|\mu\|$ measured & $\sqrt{c} / \|\mu\|$ \\
    \midrule
    \texttt{llama\_creativity} & $0.9998$ & $20.80$ & $19.16$ & $1.085$ \\
    \texttt{llama\_random} & $0.9996$ & $18.31$ & $19.16$ & $0.956$ \\
    \texttt{llama\_random\_agg\_matched} & $0.9998$ & $\text{c}<0$ & $19.16$ & n/a$^\dagger$ \\
    \texttt{qwen\_happy} & $0.9997$ & $217.31$ & $217.95$ & $0.997$ \\
    \texttt{qwen\_random} & $1.0000$ & $217.93$ & $217.95$ & $1.000$ \\
    \texttt{qwen\_random\_agg\_matched} & $0.9999$ & $217.54$ & $217.95$ & $0.998$ \\
    \texttt{qwen\_style} & $0.9947$ & $213.01$ & $217.95$ & $0.977$ \\
    \bottomrule
  \end{tabular}
  \\[2pt]
  {\footnotesize
  $^\dagger$ For aggregate-matched random on Llama, the per-layer vectors are
  rescaled large enough that even $s = 0.5$ is past the quadratic's turning
  point; the local quadratic fit is excellent (R$^2 = 0.9998$) but
  extrapolates to a negative intercept.}
\end{table}

Six of seven runs land within $\pm 10\%$ of the diagonal. \texttt{qwen\_random}
matches to four decimal places. The fit qualities are R$^2 \ge 0.9947$,
meaning the quadratic shape is essentially exact within the measured range.
The interpretation is direct: \emph{the parabola-on-log-log shape that was
visible in our $V(s)$ plots is the quadratic structure of $\|m(s)\|^2$,
and that quadratic bottoms out at $\|\mu\|^2$, which we recover}.

The second prediction --- that $V \cdot \|m\|^2 \to (\operatorname{tr}\Sigma_x - \hat{m}^\top \Sigma_x \hat{m})/2$
in the asymptotic regime --- is also borne out, but only on Llama where
the asymptotic regime is reachable. At $c = 8$ on Llama with creativity,
$V_{\mathrm{pred}}/V_{\mathrm{measured}} = 0.97$; on the aggregate-matched
random, $0.96$. On Qwen the same ratio plateaus around $1.5$ at all scales,
because Qwen's residual stream is large enough that even $c = 8$ does not
reach the asymptotic regime --- the same finding as Claim 8, restated
through the second-order lens.

The headline lesson is that the gap between the global $O(1/\|s\|)$ bound
and the observed $V(s)$ shape is not a failure of the theory. It is the
theory: \emph{at the scales we probe, we are in the crossover region of a
quadratic denominator with a $\|\mu\|^2$ floor}, and the theory predicts
exactly that shape. The slope $-1$ asymptote of Claim 8 is recovered only
once $s \gg \|\mu\|/\|v_{\mathrm{eff}}\|$.

\subsection{Threats to Validity and Meta-Findings}
\label{sec:walk-threats}

Three non-theoretical findings shaped the final form of the experiments
and merit explicit note.

\paragraph{A silent batch-indexing bug contaminated the first run.}
In HuggingFace Transformers $\ge 4.40$, the \texttt{forward} of
Qwen2/Llama3 decoder layers returns a \emph{bare} hidden-state tensor
$[B, T, d]$, not a $1$-tuple $(h,)$ as in GPT-2. An intervention helper
that did \texttt{layer.output[0][:] = layer.output[0] + s}, intending to
unpack the tuple, instead indexed into the batch dimension and steered
only the zeroth prompt of every batch. The bug was silent: no error, no
KL-gate failure (prompt $0$ did shift), CPU unit tests on \texttt{tiny-gpt2}
all passed (because GPT-2 still returns a tuple). It surfaced when a
per-layer diagnostic emitted a shape mismatch. The fix adds an
architecture-dispatching helper and shape assertions at every
capture-materialization point; a regression test verifies that all batch
elements see the intervention. All numbers reported here are post-fix;
contaminated runs are preserved at \texttt{discarded/} for audit.

\paragraph{The \texttt{welford-torch} library is unsafe in our regime.}
Our streaming covariance accumulators initially used the public
\texttt{welford-torch} library. Cross-checking against \texttt{numpy.cov}
on synthetic data with $\mu \approx 1000$ and $\sigma \approx 0.01$
revealed >$100\%$ relative error in the trace of the covariance, due to
catastrophic cancellation in a two-step delta formulation. We
hand-rolled a Chan--Golub--LeVeque batched Welford merge that stays
stable in float32. A reproducer is documented at
\texttt{docs/upstream\_bug\_reports/welford\_torch\_big\_mean\_drift/}.

\paragraph{The training pool affects the reduction gate.}
The Claim 7 failure for Qwen + happy is driven by a small but systematic
anti-alignment between the training-pool diff direction and FineWeb-Edu's
residual mean. This is not a flaw in the theory --- the gate itself is
algebraic --- but a precondition that practitioners must check before
trusting downstream bounds. We recommend reporting
$\cos(\mu, \sum_\ell s_\ell)$ as a routine diagnostic in any future work
applying these bounds.

\subsection{What to Take Away}
\label{sec:walk-takeaway}

The empirical picture matches the theory, with two important nuances.

\paragraph{Exact bounds always hold.}
Propositions~\ref{prop:walk-pole}, \ref{prop:walk-pair} (where applicable),
and Corollaries~\ref{cor:walk-V}, \ref{cor:walk-pair}, \ref{cor:walk-diam}
are strict inequalities derived from a single chordal-contraction lemma.
They hold in every run. Verifying them is a smoke test for the pipeline.

\paragraph{Asymptotic predictions hold when the scales are large enough.}
Proposition~\ref{prop:walk-local} (the local trace bound) and Claim 8 (the
$1/\|s\|$ scaling law) are correct in direction always, but their
quantitative tightness depends on whether $\|s_{\mathrm{eff}}\|$ is large
relative to the natural residual-stream scale. On Llama-3-8B at scales
the model can tolerate, the slope-$-1$ behaviour manifests cleanly. On
Qwen2.5-1.5B it never does, because Qwen's residual stream is too large
to be dominated by any tested steering magnitude.

\paragraph{Aggregate magnitude, not direction, drives the scaling.}
The single most important empirical finding is that random directions
matched in \emph{aggregate} magnitude reproduce the scaling slope of
trained vectors (Llama: $-1.85$ vs $-1.92$; Qwen as well as can be told
without an asymptotic regime). The scaling law is a statement about how
much steering you are doing, not what direction you are steering in.
Future empirical work on these bounds should report results against an
aggregate-matched random baseline rather than a per-layer-matched one.

\paragraph{The reduction gate is a precondition, not a theorem.}
The condition $\|m + s\| > \|m\|$ is algebraic; it can fail empirically
when a steering vector is mildly anti-aligned with the evaluation
distribution's residual mean. We saw this on Qwen + happy and on Qwen
single-layer interventions at L17 and L25. Practitioners applying these
bounds to a new $(\text{model}, \text{vector}, \text{dataset})$ triple
should verify the gate empirically before trusting downstream numbers.

\paragraph{The parabola on log--log has a name.}
The visible curvature of $V(s)$ is not a failure of the slope-$-1$
prediction; it is the correct second-order Taylor shape, with
$V(s) \propto 1/\|m(s)\|^2$ and $\|m(s)\|^2$ a quadratic in $s$ with
floor $\|\mu\|^2$. Fitting the quadratic recovers $\sqrt{c} = \|\mu\|$
within $1$--$9\%$ across every run with a valid local fit. The slope-$-1$
asymptote is what one sees only deep into the regime
$\|s_{\mathrm{eff}}\| \gg \|\mu\|$.

\paragraph{Open: the eigenalignment claim.}
Claim 9's prediction --- that the smallest-eigenvalue eigenvector of
$\Sigma_z$ aligns with $\hat{m}$ --- is the one prediction in the paper
that we do not currently see in the data. We do not yet know whether the
issue is statistical (noisy estimate of the smallest eigenvector with
$K = 1024$ samples in high dimension), structural (the propagated
$\hat{m}$ at the measurement layer is rotated relative to the injection
direction), or substantive. Future work that evaluates Claim 9 at the
injection layer is the natural next step.
